\cotefile{12}
\scanpage{64}
% Three vertical emphasis strokes to the left of lines 1–3.
soient des isomorphismes, avec % 1
le complément que $H^*(X,M)$ \uline{est un} % 2
$\mathbb{Q}$-\uline{\deleted{\unsure{m}}module de type fini}. \deleted{Ceci doit être % 3
valable en particulier dans les deux % 4
cas types suivants % 5

a) $X$ de type fini sur $\mathbb{Z}$ % 6

b) $X$ de type fini sur un corps % 7
alg. clos.} % 8

En particulier, pour un corps alg. clos $k$, on % 9
aura $H^i(k,M)=0$ pour $i>0$, $\forall$ motif $M$. % 10

\uline{NB} Si $X$ est de type fini (ou projectif lisse) % 11
sur un corps alg. clos $K$, on \unsure{n'a} plus de % 12
telles relations précises, comme on voit % 13
par calcul \unsure{confiant|cofini|explicite} (utilisant les conjectures % 14
de Tate) pour $K$ clôture algébrique % 15
d'un corps fini ---- % 16

Le th de Tate permet de reformuler et % 17
se préciser en énonçant : % 18
\[
 A^i(X)\otimes_{\mathbb{Z}}\overset{\inserted{\nearrow\,\text{pas}\,\mathbb{Q}_{\ell}}}{\mathbb{Q}}
 \simeq\Gamma_k\bigl(R^i f_*(\mathbb{Q}_X(0))\bigr) % 19
\]
% archive sheet number: 63
